\section{High-dimensional orthogonality}
\label{sec:orthogonality}

\begin{lemma}[Projection variance and concentration on the sphere]
\label{lem:orthogonality}
Let $u\in\R^d$ be uniformly distributed on the sphere $S^{d-1}$ and let
$e\in\R^d$ be any deterministic unit vector. Then
\begin{align}\label{eq:orthogonality_variance}
   \E\inner{e}{u}^2 \;=\; \frac1d,
\end{align}
and for every $t\in[0,c'\sqrt d\,]$ with $c'\in(0,1)$ an absolute
constant,
\begin{align}\label{eq:orthogonality_tail}
   \Pr\!\left[\, \abs{\inner{e}{u}} \;>\; \frac{t}{\sqrt d} \,\right]
   \;\le\; 2\exp(-c\,t^2),
\end{align}
for an absolute constant $c>0$. Moreover $\inner{e}{u}$ obeys a matching
small-ball (anti-concentration) lower bound: for $t$ in a bounded range,
$\Pr[\inner{e}{u}>t/\sqrt d]\ge c''\exp(-C''t^2)$ for absolute constants
$c'',C''>0$.
\end{lemma}

\begin{proof}[Proof of \Cref{lem:orthogonality}]
The proof is three named facts about the uniform measure on $S^{d-1}$, two
of them direct from \cite{vershynin2018}.

\smallskip\noindent\textbf{Variance \Cref{eq:orthogonality_variance}.} By
rotational invariance of the uniform measure, $\E[uu^\top]$ commutes with
every orthogonal matrix, hence is a scalar multiple of the identity,
$\E[uu^\top]=\lambda I$; taking traces, $1=\E\norm{u}_2^2=\lambda d$, so
$\lambda=1/d$ and $\E\inner{e}{u}^2=e^\top\E[uu^\top]e=\norm{e}_2^2/d=1/d$.

\smallskip\noindent\textbf{Upper tail \Cref{eq:orthogonality_tail}.} The
linear form $\inner{e}{u}$ of a uniform point on $S^{d-1}$ is sub-Gaussian
with variance proxy of order $1/d$: \cite{vershynin2018}, Theorem~3.4.6
(concentration of Lipschitz functions on the sphere; $u\mapsto\inner{e}{u}$
is $1$-Lipschitz) gives $\Pr[\abs{\inner{e}{u}}>r]\le2\exp(-c d r^2)$ for an
absolute $c>0$; substituting $r=t/\sqrt d$ yields
$\Pr[\abs{\inner{e}{u}}>t/\sqrt d]\le2\exp(-c t^2)$, valid for
$t\in[0,c'\sqrt d]$.

\smallskip\noindent\textbf{Small-ball lower bound.} The exact density of
$\inner{e}{u}$ is $f(r)=c_d(1-r^2)^{(d-3)/2}$ on $[-1,1]$ with
$c_d=\Theta(\sqrt d)$ (the projection of the uniform sphere measure; see
\cite{vershynin2018}, Section~3.4). For $r=t/\sqrt d$ with $t$ in a bounded
range, $(1-r^2)^{(d-3)/2}=\exp\bigl(\tfrac{d-3}{2}\log(1-t^2/d)\bigr)\ge\exp(-C''t^2)$
for large $d$, so integrating $f$ over $[t/\sqrt d,\,2t/\sqrt d]$ (an
interval of length $t/\sqrt d$) gives
$\Pr[\inner{e}{u}>t/\sqrt d]\ge c''\exp(-C''t^2)$ for absolute constants
$c'',C''>0$. This is the spreading estimate feeding the anti-concentration
argument of \Cref{lem:incorrect_max_lower}. This completes the proof.
\end{proof}

\begin{remark}[The structural role of $1/d$, and that we do not use L\'evy for $1/\sqrt T$]
\label{rem:orthogonality_structural}
\Cref{eq:orthogonality_variance} is the quantitative form of ``random
vectors are nearly orthogonal in high dimension''. It is the sole source
of the $1/d$ factor in the per-step noise variance proxy
$\E[M_t^2\mid\Fcal_{t-1}]\le\sigma^2/d$ used by \Cref{lem:azuma}, and
hence of the $1/\sqrt d$ half of the noise scale. The complementary
$1/\sqrt{T}$ half comes from the softmax quadratic variation
$S_T\le e^{2S}/T$ (\Cref{lem:quad_variation}), \emph{not} from a
sphere concentration-of-measure (L\'evy) argument: the high-dimensional
spreading here governs single-direction projection variance, while the
horizon dependence is purely the running-average weight structure.
\end{remark}
